\chapter*{Περίληψη}\addcontentsline{toc}{chapter}{Περίληψη}
\thispagestyle{plain}

Η παρούσα διπλωματική εργασία πραγματεύεται την υλοποίηση ενός
πειραματικού Χώρου Δεδομένων (\textlatin{Data Space}) με
χρήση των ανοιχτών τεχνολογιών \textlatin{FIWARE} και του πλαισίου εμπιστοσύνης \textlatin{i4Trust}/\textlatin{iSHARE}.
Βασικός σκοπός της εργασίας είναι η υλοποιήση του μηχανισμού
που εξασφαλίζει την ψηφιακή κυριαρχία (\textlatin{digital sovereignty}) κατά την κοινή χρήση και ανταλλαγή δεδομένων μεταξύ διαφορετικών οργανισμών.

Στο πλαίσιο της εργασίας υλοποιήθηκε ένα
πλήρως λειτουργικό πρωτότυπο \textlatin{Data Space} που αποδεικνύει τον κύκλο
εκχώρησης πρόσβασης σε δεδομένα μεταξύ δύο ανεξάρτητων οντοτήτων. Η
αρχιτεκτονική οργανώνεται σε τρία λογικά επίπεδα: το επίπεδο επιβολής πολιτικής
(\textlatin{Kong API Gateway} και \textlatin{DSBA-PDP}), το επίπεδο εμπιστοσύνης
και ταυτότητας (\textlatin{iSHARE Satellite} και \textlatin{Keyrock Identity
Manager}) και το επίπεδο δεδομένων (\textlatin{Orion-LD Context Broker}).

Ο έλεγχος πρόσβασης βασίζεται στο πρωτόκολλο \textlatin{iSHARE}, που
αξιοποιεί ψηφιακά πιστοποιητικά \textlatin{X.509} και \textlatin{JSON Web
Tokens (JWT)} για την ταυτοποίηση των συμμετεχόντων, καθώς και μηχανισμό
εκχώρησης εξουσιοδότησης (\textlatin{delegation evidence}) για τον ορισμό και την
επιβολή πολιτικών πρόσβασης (\textlatin{access policies}). Τα δεδομένα αποθηκεύονται και ανταλλάσσονται
σε μορφή \textlatin{NGSI-LD}, σύμφωνα με τα \textlatin{Smart Data Models} και το πρότυπο \textlatin{ETSI}.

Το σύστημα τροφοδοτήθηκε με πραγματικά δεδομένα από το \textlatin{Plegma
Dataset}, μια συλλογή μετρήσεων ζήτησης ενέργειας νοικοκυριών, περιβαλλοντικών
παραμέτρων και μεταδεδομένων κτιρίων. Η επίδειξη του κύκλου
εκχώρησης πρόσβασης αποδεικνύει επιτυχώς την πλήρη πρόσβαση του παρόχου στα δικά του δεδομένα, 
την απόρριψη μη εξουσιοδοτημένων αιτημάτων, και την
ελεγχόμενη πρόσβαση του καταναλωτή αποκλειστικά στους πόρους για τους οποίους
ο πάροχος έχει εκχωρήσει ρητή άδεια χρήσης.

\vspace{1cm}
\flushleft \textbf{Λέξεις Κλειδιά:} \textlatin{Data Spaces}, \textlatin{FIWARE},
\textlatin{i4Trust}, \textlatin{iSHARE}, \textlatin{NGSI-LD}, Έλεγχος Πρόσβασης,
\textlatin{JSON Web Tokens}, Εκχώρηση Εξουσιοδότησης, Ψηφιακή Κυριαρχία
\blankpage